\section{The positive quadratic conductor-eight rows from the existing unit certificates}
\label{sec:positive-quadratic8-strict}

\paragraph{Scope.}
This section concerns the actual finite character $\chi_8$ of
$\mathbb Q(\sqrt2)$, with no modulus-five factor.
It uses the already proved class and unit inputs of
AC-G142-class and AC-G147-class; no new analytic or class-number
certificate is required.
The new calculation is the quadratic, rather than quartic,
projector on those same exact global units.
We explicitly allow the residual character to have a kernel on
the chosen larger prime-to-$269$ cover.
The negative character $\chi_{-8}$ is not included in this theorem.

Put $p=269$, $K=\mathbb Q(\omega)$ and $t=\sqrt[3]{p}$.
Let $\chi:G_K\to\mu_3$ describe $K(t)/K$, and let
$\xi=\chi_8|_{G_K}$.
Use the actual cyclotomic character $c_p$, its Teichmuller
part $\omega_p$, and the quadratic character
$\tau_p=\omega_p^{134}$ of $\mathbb Q(\sqrt p)$.
For a discrete module $D$ set
\[
 \begin{gathered}
 S_{\rm str}(K,D)=
 \ker\bigl(H^1_{\rm cts}(G_K,D)\longrightarrow\mathcal L(D)\bigr),\\
 \mathcal L(D)=
 \prod_{v\mid p}H^1(G_{K_v},D)\ \times\
 \prod_{v\nmid p,\infty}H^1(I_v,D).
 \end{gathered}
 \tag{Q8P.1}
\]

\begin{theorem}[AC-G149: the actual positive quadratic-eight strict rows]
\label{AC-G149}
Let $E/\mathbb Q_p$ be finite and $O$ its original ring of integers,
containing the values of $\chi$.
For $j\in\{1,2\}$, $m\in\mathbb Z$, $e\geq1$, and any continuous
$\kappa:G_K\to1+\mathfrak m_O$ unramified away from $p$, put
\[
 D=(E/O)(\chi^{-j}\xi\omega_p^{-m}\kappa^{-1}c_p^{-e}).
 \tag{Q8P.2}
\]
If $e+m\equiv267$ or $133\bmod268$, then
\[
                         S_{\rm str}(K,D)=0 .
 \tag{Q8P.3}
\]
The corresponding residual strict group is zero over every finite
field containing the character values.
No absolute ramification restriction is imposed on $O$.
For $m=0,\kappa=1$, these are the two original finite labels
$\chi^{-j}\chi_8$ of the actual AC-K25 hosts.
The residual quadratic character $\tau_p$ in the second case
does not become a new finite host label.
\end{theorem}

\begin{proof}
\proofstep{1. Existing covers and the inherited quadratic class inputs.}
Let $\phi$ be the primitive odd conductor-sixteen character with
$\phi(13)=i$, $\phi(5)=-i$, $\phi^2=\chi_8$ and
$i\bmod p=82$.
Choose the following two already certified quartic covers:
\[
 \begin{gathered}
 F_{267}=\mathbb Q(w),\qquad w^4+4w^2+2=0,\\
 F_{133}=\mathbb Q(z),\qquad z^4-1076z^2+4842=0,\\
 \begin{array}{c|c|c}
 r=e+m\bmod268&\theta_r&\theta_r^2\\ \hline
 267&\phi&\chi_8\\
 133&\phi\omega_p^{67}&\chi_8\tau_p .
 \end{array}
 \end{gathered}
 \tag{Q8P.4}
\]
The second field and its character orientation are exactly
$F_1$ of AC-G147-class, not a field identified only by its discriminant.
Put $L_r=F_r(t)$ and $M_r=KL_r$.
In both cases
\[
 \begin{gathered}
 \operatorname{Gal}(M_r/\mathbb Q)=S_3\times C_4,\\
 \Delta_r=\operatorname{Gal}(M_r/K)=C_3\times C_4,\qquad
 \operatorname{Cl}(L_r)[p]=0.
 \end{gathered}
 \tag{Q8P.5}
\]
The class assertions are respectively AC-G142-class and AC-G147-class.
Each $M_r$ has degree twenty-four over $\mathbb Q$, so
$\mu_p\not\subset M_r$.
Let $d$ generate $C_3$ by $d(t)=\omega t$, and let $b$
generate $C_4$, fixing $K,t$ and having $\theta_r(b)=i$.

The unique quadratic subfields of $F_{267}$ and $F_{133}$
are $\mathbb Q(\sqrt2)$ and $\mathbb Q(\sqrt{538})$.
Thus
\[
 [L_{267}:\mathbb Q(t,\sqrt2)]=2,\qquad
 [L_{133}:\mathbb Q(t,\sqrt{538})]=2 .
 \tag{Q8P.6}
\]
For each inclusion, class extension followed by norm is multiplication
by two on the smaller class group.
The extension map is therefore injective on its $p$-primary part,
since two is a unit at $p$.
Consequently both smaller class groups have zero $p$-part too.
This is an inheritance from existing class certificates, not a
new numerical class-group computation.
We shall retain the larger covers because their exact unit witnesses
are already available.

\proofstep{2. Kummer descent does not require a faithful cover character.}
Write $k=O/\mathfrak m_O$.
The residual character of (Q8P.2) and its Kummer quotient are
\[
 \begin{aligned}
 a&=\bar\chi^{-j}\bar\xi\,\bar c_p^{-(e+m)},\\
 \eta:=a/\bar c_p&=
 \begin{cases}
 \bar\chi^{-j}\bar\xi,&r=267,\\
 \bar\chi^{-j}\bar\xi\,\bar\tau_p,&r=133.
 \end{cases}
 \end{aligned}
 \tag{Q8P.7}
\]
In the corresponding cover this is
$\eta=\bar\chi^{-j}\bar\theta_r^{\,2}$.
Its kernel on $\Delta_r$ is the subgroup $\langle b^2\rangle$.
Its minimal splitting field is $K(t,\sqrt2)$ or $K(t,\sqrt{538})$,
respectively, of index two in $M_r$.
No faithfulness is needed in the following descent.
The coefficient restricts to $k(1)$ on $G_{M_r}$, and
$|\Delta_r|=12$ is invertible in $k$.
Higher cohomology of this finite quotient vanishes by averaging.
Inflation--restriction therefore identifies the global degree-one
cohomology with the appropriate twisted invariants over $M_r$.
The same argument applies to every local decomposition or inertia
quotient; all their orders divide twelve.
These identifications commute with restriction.

In particular the residual strict kernel is
\[
 e_{\eta^{-1}}\ker\left(
 V_{M,S}\otimes_{\mathbb F_p}k
       \longrightarrow\mathcal P_M\otimes_{\mathbb F_p}k\right),
 \tag{Q8P.8}
\]
where
\[
 \begin{gathered}
 V_{M,S}=\{x\in M^\times:p\mid v(x)\text{ for }v\nmid p\}/M^{\times p},\\
 \mathcal P_M=\bigoplus_{v\mid p}M_v^\times/M_v^{\times p}.
 \end{gathered}
 \tag{Q8P.9}
\]
The inverse character occurs because Kummer cohomology has the
additional twist $\eta$.
It has cubic eigenvalue $\bar\chi^j$ and $b$-eigenvalue $-1$.
The extra subgroup $\langle b^2\rangle$ is taken into account
by these invariants; it is not discarded.

Full local zero at $p$ forces the Kummer valuations there to
be divisible by $p$, as do the away inertia conditions at
every other finite place.
The kernel lies in the all-primes valuation group $V_M$, with
the exact sequence
\[
 0\longrightarrow\mathcal O_M^\times/\mathcal O_M^{\times p}
 \longrightarrow V_M\longrightarrow\operatorname{Cl}(M)[p]
 \longrightarrow0.
 \tag{Q8P.10}
\]
For the reflection $s$ fixing $L_r$, norm followed by extension
on $\operatorname{Cl}(M_r)[p]$ is $1+s$.
The norm is zero by (Q8P.5), so $s=-1$ there.
Every conjugate reflection acts the same way, and their products
generate the cubic subgroup.
Hence $C_3$ acts trivially on $\operatorname{Cl}(M_r)[p]$.
All nontrivial cubic class sectors vanish.
This does not claim that the entire class group of $M_r$ vanishes.

\proofstep{3. The quadratic unit sector has the required dimension.}
In either cover, the twelve complex places over the chosen complex
place of $K$ form a regular $\Delta_r$-set.
The equivariant Dirichlet unit character is
$\mathbb Q[\Delta_r]-\mathbb Q$.
Because $p\nmid12$ and $\mu_p(M_r)=1$, reduction modulo $p$
preserves these multiplicities.
Let $W_-$ be the nontrivial cubic part of the $b=-1$ eigenspace
in the units modulo $p$.
It has dimension two over $\mathbb F_p$ and is simple under $C_3$,
since $X^2+X+1$ is irreducible modulo $269$.
The direct sum of the local Kummer maps is $\Delta_r$-equivariant.
It suffices to find a single nonzero local value of a global
class in $W_-$.
After scalar extension this proves injectivity on both cubic lines.
This argument applies also though the finite quotient character
being selected is not faithful.

\proofstep{4. The first residue and both unramified quartic completions.}
Use the exact unit number four of AC-G142.
Its rational tensor vector and monic constant-one integral relation
are already certified there; it is an actual global unit, with no
fundamental-unit completeness assumption.
For this cover, there are two places of $M$ above $p$, each with
ramification degree three and residue degree four.
The element $b$ exchanges them.
At a chosen place, $t$ is a uniformizer and the residue field is
$\mathbb F_p[w]/(w^4+4w^2+2)$.
The exact reductions are
\[
 \begin{gathered}
 U=A+tB\pmod{t^2},\\
 A=26+232w^2,\qquad
 B=143+123w+169w^2+39w^3,\\
 R=3B/A=263+135w+266w^2+252w^3.
 \end{gathered}
 \tag{Q8P.11}
\]
Here $A$ is nonzero in this degree-four field.
Form the global norm-corrected unit
$V=U^3/\prod_{a=0}^2d^aU$.
Its local expansion is $1+tR\bmod t^2$.
Apply the global quadratic projector
\[
                  P_-=1-b+b^2-b^3=(1-b)(1+b^2).
 \tag{Q8P.12}
\]
In multiplicative language this means the product of the indicated
conjugate units with integer exponents $1,-1,1,-1$.
Since $b(w)=-3w-w^3$ and $b(t)=t$, its linear coefficient is
\[
 \sum_{a=0}^3(-1)^aR(b^a w)=245+257w^2\ne0.
 \tag{Q8P.13}
\]
This exact finite-polynomial calculation is independently checked
in the accompanying new projector script.
Explicitly $b(w^2)=-4-w^2$, while $1+b^2$ kills the odd powers
of $w$, which gives the displayed coefficient directly.
The product is formed globally before local evaluation.
We use both local completions in (Q8P.9), and do not assign a
$K_p$-linear action of $b$ to just one of them.
The nonzero value at the displayed place gives a nonzero vector
in the full direct-sum target.
A local $p$th power unit has zero coefficient of $t$ modulo $t^2$:
the unramified coefficient field supplies residue lifts, and
the linear binomial term is divisible by $pt=t^4$.
Thus (Q8P.13) proves nonzero localization on $W_-$ for $r=267$.

\proofstep{5. The second residue and its unique ramified completion.}
Use unit number five in the first field of AC-G147.
Its full numerator vector, common denominator $69p$ and monic
constant-one relation are fixed by
\path{selected_units_exact.json} in that proof.
The exact rational change of basis there retains every denominator
before reduction, and proves all local coefficients integral.
Here $M$ has one prime above $p$, with $e=12,f=2$.
Use $u=t/z$ and the exact Eisenstein equation
\[
                   18^3u^{12}+p(12\cdot18-64p)u^6+p=0.
 \tag{Q8P.14}
\]
The entire local group $\Delta_r$ stabilizes this one place.
Modulo $u^{12}$ its quartic action is
$b(u)=187u+262u^7$, and its cubic action is $d(u)=\omega u$.
The normalized logarithm of the certified unit, in degrees
zero through eleven, is
\[
                 (0,140,88,127,0,231,12,233,60,227,192,211).
 \tag{Q8P.15}
\]
Apply the global cubic norm correction and then (Q8P.12).
On the logarithm, the first operator kills degrees divisible by
three and multiplies all other degrees by three.
The second uses the actual nonlinear $b$-substitution.
The result is
\[
                     249u^2+182u^8+152u^{10}\ne0
                         \pmod{u^{12}}.
 \tag{Q8P.16}
\]
The new finite verifier checks this expression and its eigenvalue
$b(P_-L)=-P_-L$ directly from the full certified unit reduction.
In particular it retains the $u^7$ part of the automorphism.
For a short direct check, the even part after cubic correction is
$264u^2+180u^8+38u^{10}$.
Here $b(u^2)=-u^2+72u^8$, $b(u^8)=u^8$ and
$b(u^{10})=-u^{10}$ modulo $u^{12}$.
On this even part $P_-=2(1-b)$, giving exactly (Q8P.16).
Since $v_u(p)=12<p-1$, logarithm and exponential identify
principal units with the maximal ideal.
A $p$th power has logarithm in its thirteenth power.
The nonzero terms in (Q8P.16) therefore give a nonzero local
Kummer class in $W_-$.
By Step3, localization is injective on both cubic eigenlines
also in this second residue.
Together with (Q8P.8)--(Q8P.10) and the vanished class sector,
this proves residual strict zero in both cases.

\proofstep{6. The original coefficient uniformizer and all-primary conclusion.}
Let $\varpi$ be an original uniformizer of $O$.
The module $D$ is $\varpi$-divisible and $D[\varpi]\simeq k(a)$.
At $p$, the quadratic conductor-eight character is unramified.
The residual cubic inertia has order three and cannot cancel
the residual cyclotomic factor, whose order divides $268$.
The reduction of $\kappa$ is trivial.
Thus $D^{G_{K_p}}=0$.
Away from $p$, the cyclotomic, Teichmuller and $\kappa$ factors
are unramified and the finite inertia has order dividing six.
Its invariants are either zero or the entire divisible module $D$.
For every local group used in (Q8P.1),
$D^J/\varpi D^J=0$.
The long exact sequence of multiplication by $\varpi$
therefore gives local injectivity
$H^1(J,D[\varpi])\to H^1(J,D)$.
Every global strict class killed by $\varpi$ lifts to a residual
strict class, which is zero by the preceding steps.
Finally a continuous cocycle from compact $G_K$ to discrete $E/O$
has finite image and is killed by a power of $\varpi$.
A nonzero strict class would produce a nonzero one killed by
$\varpi$, a contradiction.
This proves the theorem on the original $O$ without replacing it.
\end{proof}

\paragraph{Actual labels, count, and retained scope.}
The actual AC-K25 hosts contain $\sqrt2$ in their two-power
cyclotomic factor and contain $K(t)$.
Thus the two labels $\chi^{-j}\chi_8$ occur there.
They are distinct from the conductor-four, conductor-sixteen
and pure-cubic labels in AC-G146/148.
For $m=0,\kappa=1$, their two positive Tate progressions combine
to $e=134n-1$ and have count
$2\lfloor T/134\rfloor$ for $1\leq e<T$.
Neither the larger quartic covers nor $\tau_p$ add a finite
character to the original coefficient in this statement.
Transfer to the integral-index window still requires the
retained local comparison cost; no individual index is set to zero.

\paragraph{Exact evidence and source scope.}
The new exact calculation is
\path{verify_quadratic8_projectors.py} with its complete JSON output.
It checks (Q8P.13) and (Q8P.16) from the already certified actual
unit reductions and automorphism formulas.
The all-denominator conversions and monic unit relations remain
the pinned AC-G142 and AC-G147 evidence, not a new unit-basis assertion.
Class inputs, equivariant Dirichlet units, restriction/corestriction
and Kummer descent retain their earlier complete applied scopes.
The Kummer sequence is the retained Stacks tag~03PK input.
The arithmetic finite-cover dictionary is the retained SGA~1,
Expos\'e~V, \S\S4 and 7 dictionary.
No new analytic interval, class-number computation, primary-source
proof, Lean execution, mixed realization or all-prime estimate is claimed.
